\chapter{Jaring Makanan}\label{ch:g5:food-webs}

\Cref{ch:g3:food-chains} sudah memperingatkan kita bahwa rantainya
yang rapi itu hanyalah utas tunggal yang ditarik dari sebuah jala:
elangnya juga menyambar tikus, kadalnya juga memakan lalat, rumputnya
memberi makan kelinci pula. Tahun ini kita berhenti menarik benang
lalu memandang seluruh jalanya --- \omterm{def:g5:food-webs:web}{jaring makanan} --- dan memandang
apa yang menjaganya tetap seimbang.

\section{Dari rantai menjadi jaring}

\begin{definition}[Jaring makanan]\label{def:g5:food-webs:web}
Yang disebut \emph{jaring makanan}\index{jaring makanan} sebuah
\omterm{def:g2:where-animals-live:habitat}{habitat} adalah jaringan \emph{semua} \omterm{def:g3:food-chains:chain}{rantai makanannya} sekaligus:
setiap \omterm{def:g5:biodiversity:species}{spesies} digambar satu kali, dengan sebuah anak panah (yang
masih berarti \emph{\omterm{def:g3:food-chains:chain}{dimakan oleh}},
\cref{rem:g3:food-chains:arrow}) untuk setiap hubungan santapan.
Kebanyakan \omterm{def:g1:animals-around-us:animal}{hewan} memakan beberapa \omterm{def:g5:biodiversity:species}{spesies} dan \omterm{def:g3:food-chains:chain}{dimakan oleh} beberapa
\omterm{def:g5:biodiversity:species}{spesies} --- sehingga rantainya bersilangan, dan gambarnya menjadi
sebuah jaring.
\end{definition}

\begin{omfigure}
\begin{tikzpicture}[scale=0.98]
  % nodes
  \node[draw, thick, omMeth, fill=omMeth!8, rounded corners=3pt,
    font=\small] (grass) at (0,0) {rumput};
  \node[draw, thick, omMeth, fill=omMeth!8, rounded corners=3pt,
    font=\small] (clover) at (2.6,0) {semanggi};
  \node[draw, thick, omMeth, fill=omMeth!8, rounded corners=3pt,
    font=\small] (oak) at (5.2,0) {daun pohon ek};
  \node[draw, thick, omExo, fill=omExo!8, rounded corners=3pt,
    font=\small] (hopper) at (0,1.7) {belalang};
  \node[draw, thick, omExo, fill=omExo!8, rounded corners=3pt,
    font=\small] (rabbit) at (2.6,1.7) {kelinci};
  \node[draw, thick, omExo, fill=omExo!8, rounded corners=3pt,
    font=\small] (caterpillar) at (5.2,1.7) {ulat};
  \node[draw, thick, omProp, fill=omProp!8, rounded corners=3pt,
    font=\small] (lizard) at (0,3.4) {kadal};
  \node[draw, thick, omProp, fill=omProp!8, rounded corners=3pt,
    font=\small] (tit) at (5.2,3.4) {gelatik biru};
  \node[draw, thick, omThm, fill=omThm!8, rounded corners=3pt,
    font=\small] (fox) at (2.6,4.4) {rubah};
  \node[draw, thick, omThm, fill=omThm!8, rounded corners=3pt,
    font=\small] (buzzard) at (0.2,5.2) {elang};
  \node[draw, thick, omThm, fill=omThm!8, rounded corners=3pt,
    font=\small] (hawk) at (5.0,5.2) {alap-alap};
  % arrows
  \draw[->, thick, omDef] (grass) -- (hopper);
  \draw[->, thick, omDef] (grass) -- (rabbit);
  \draw[->, thick, omDef] (clover) -- (rabbit);
  \draw[->, thick, omDef] (oak) -- (caterpillar);
  \draw[->, thick, omDef] (hopper) -- (lizard);
  \draw[->, thick, omDef] (caterpillar) -- (tit);
  \draw[->, thick, omDef] (rabbit) -- (fox);
  \draw[->, thick, omDef] (lizard) -- (buzzard);
  \draw[->, thick, omDef] (hopper) -- (tit);
  \draw[->, thick, omDef] (tit) -- (hawk);
  \draw[->, thick, omDef] (rabbit) -- (buzzard);
\end{tikzpicture}

{\small Sebuah \omterm{def:g5:food-webs:web}{jaring makanan} saku dari satu tepi padang rumput.
Setiap rantai lurus pada \cref{ch:g3:food-chains} ada di dalamnya ---
bersilangan dengan semua yang lain.}
\end{omfigure}

\begin{example}[Membaca jaringnya]\label{ex:g5:food-webs:reading}
Dua cara membaca gambarnya. Ikutilah anak panah kelincinya: ia memakan
rumput \emph{dan} semanggi, dan dimakan rubah \emph{dan} elang ---
empat hubungan, sedangkan sebuah rantai hanya memperlihatkan dua. Lalu
hitunglah anak panah makanan elangnya: kadal dan kelinci --- hilang
satu, yang lain tetap ada. Jaringnya mencatat persis pilihan cadangan
yang dipuji \cref{prop:g5:biodiversity:web}.
\end{example}

\section{Peran di dalam jaringnya}

\begin{definition}[Produsen dan konsumen]\label{def:g5:food-webs:roles}
\omterm{def:g5:biodiversity:species}{Spesies} di dalam jaringnya terkelompok menurut peran:
\begin{itemize}
  \item \emph{produsen}\index{produsen}: \omterm{def:g1:plants-around-us:plant}{tumbuhan}, yang membuat
        materinya sendiri
        (\cref{prop:g4:what-plants-need:make}) --- setiap jalur anak
        panah bermula pada salah satunya
        (\cref{prop:g3:food-chains:start});
  \item \emph{konsumen}\index{konsumen}: \omterm{def:g1:animals-around-us:animal}{hewan}, yang hidup dari yang
        lain --- \omterm{def:g2:what-animals-eat:kinds}{herbivora} memakan produsen secara langsung,
        \omterm{def:g2:what-animals-eat:kinds}{karnivora} memakan konsumen lain, \omterm{def:g2:what-animals-eat:kinds}{omnivora} memakan keduanya
        (\cref{def:g2:what-animals-eat:kinds});
  \item dan, yang bekerja di balik bayang-bayang,
        \emph{pengurai}\index{pengurai} --- para pelapuk pada
        \cref{ex:g2:caring-for-nature:compost}, yang makan dari sisa
        yang mati lalu mengembalikan zatnya ke dalam tanah. Tahun
        berikutnya memberi mereka bab tersendiri.
\end{itemize}
\end{definition}

\begin{example}[Mengelompokkan gambarnya]\label{ex:g5:food-webs:sorting}
Di dalam jaring sakunya: \omterm{def:g5:food-webs:roles}{produsen} --- rumput, semanggi, pohon ek;
\omterm{def:g5:food-webs:roles}{konsumen} \omterm{def:g2:what-animals-eat:kinds}{herbivora} --- belalang, kelinci, \omterm{ex:g2:animals-grow-up:butterfly}{ulat}; \omterm{def:g5:food-webs:roles}{konsumen} \omterm{def:g2:what-animals-eat:kinds}{karnivora} ---
kadal, gelatik biru (kebanyakan), rubah (sesungguhnya seorang
\omterm{def:g2:what-animals-eat:kinds}{omnivora}), elang, alap-alap. Tiga lantai: hijau, pemakan \omterm{def:g1:plants-around-us:plant}{tumbuhan},
pemakan daging --- dengan setiap lantai atas berdiri di atas lantai di
bawahnya.
\end{example}

\begin{proposition}[Bentuk bilangannya]\label{prop:g5:food-webs:pyramid}
Menghitung penghuni jaringnya lantai demi lantai memberi sebuah bentuk
yang tetap: \omterm{def:g5:food-webs:roles}{produsen} \emph{banyak}, \omterm{def:g2:what-animals-eat:kinds}{herbivora} \emph{lebih sedikit},
\omterm{def:g2:what-animals-eat:kinds}{karnivora} \emph{sedikit} --- sebuah padang rumput berisi jutaan
tanaman rumput memberi makan ribuan belalang, yang memberi makan
segelintir kadal, yang memberi makan sepasang elang. Tiap lantai hidup
dari lantai di bawahnya dan tidak dapat melebihi jumlah makanannya.
Itulah sebabnya pemburu besar selalu langka.
\end{proposition}
\admitted

\section{Keseimbangan, dan apa yang mengganggunya}

\begin{proposition}[Keseimbangan jaringnya]\label{prop:g5:food-webs:balance}
Kalau dibiarkan sendiri, sebuah jaring menjaga bilangannya dalam
\emph{keseimbangan}\index{keseimbangan jaring makanan} yang kasar:
kalau kelinci berlipat ganda, pemburunya makan lebih baik lalu
berlipat ganda sesudahnya, dan kelebihannya meleleh; kalau kelinci
menyusut, pemburu yang lapar beralih ke mangsa lain (anak panah
cadangannya) dan jumlah kelincinya pulih. Pembetulan pada
\cref{prop:g3:food-chains:depend}, yang berjalan di sepanjang banyak
anak panah sekaligus, memantapkan keseluruhannya.
\end{proposition}
\admitted

\begin{example}[Kekeliruan yang klasik]\label{ex:g5:food-webs:mistake}
Para petani dahulu berkampanye untuk memusnahkan \omterm{prop:g3:sorting-living-things:groups}{burung} pemangsa ---
``pencuri ayam''. Di tempat mereka berhasil, tikus dan tikus sawah,
yang terbebas dari pemburu utamanya, berlipat ganda menjadi wabah yang
memakan gandum jauh lebih banyak daripada nilai satu ayam yang sesekali
diambil elangnya. Kisah siput sang tukang kebun
(\cref{exo:g3:food-chains:9}), pada skala seluruh pedesaan:
singkirkan lantai teratas sebuah jaring dan lantai di bawahnya
meledak.
\end{example}

\begin{method}[Meramalkan sebuah gangguan]\label{met:g5:food-webs:predict}
Untuk meramalkan akibat menyingkirkan (atau menambahkan) sebuah
\omterm{def:g5:biodiversity:species}{spesies}:
\begin{enumerate}
  \item temukan spesiesnya di dalam jaring lalu daftarkan \emph{semua}
        anak panahnya, yang masuk dan yang keluar;
  \item makanannya, yang lebih sedikit dimakan, cenderung bertambah;
        pemakannya, yang menjadi lebih miskin, cenderung berkurang
        atau berganti mangsa;
  \item ikuti tiap perubahan satu langkah lagi di sepanjang jaringnya
        --- akibat langkah keduanya sering mengejutkan;
  \item periksa tetangga mana yang punya anak panah cadangan (mereka
        menyerap guncangannya) dan mana yang bersandar hanya pada
        \omterm{def:g5:biodiversity:species}{spesies} ini (mereka dalam bahaya).
\end{enumerate}
\end{method}

\section{Latihan}

\begin{exercise}[$\star$]\label{exo:g5:food-webs:1}
Apa itu \omterm{def:g5:food-webs:web}{jaring makanan}, dan mengapa rantainya bersilangan?
\end{exercise}

\begin{exercise}[$\star$]\label{exo:g5:food-webs:2}
Sebutkan ketiga peran pada \cref{def:g5:food-webs:roles}, dengan satu
contoh masing-masing dari gambarnya (kecuali \omterm{def:g5:food-webs:roles}{pengurai}).
\end{exercise}

\begin{exercise}[$\star$]\label{exo:g5:food-webs:3}
Dari gambarnya: daftarkan segala yang dimakan gelatik birunya, dan
segala yang memakannya.
\end{exercise}

\begin{exercise}[$\star$]\label{exo:g5:food-webs:4}
Tariklah dari gambarnya dua rantai lengkap yang berbagi belalangnya.
\end{exercise}

\begin{exercise}[$\star$]\label{exo:g5:food-webs:5}
Nyatakan bentuk bilangannya, lantai demi lantai. Mengapa tidak ada
lantai yang dapat melebihi jumlah lantai di bawahnya?
\end{exercise}

\begin{exercise}[$\star$]\label{exo:g5:food-webs:6}
Mengapa elang langka dan belalang tak terhitung, dalam satu kalimat?
\end{exercise}

\begin{exercise}[$\star$]\label{exo:g5:food-webs:7}
Kelinci mengalami tahun yang baik lalu berlipat ganda. Menurut
\cref{prop:g5:food-webs:balance}, apa yang terjadi berikutnya --- dan
apa yang melelehkan kelebihannya?
\end{exercise}

\begin{exercise}[$\star$]\label{exo:g5:food-webs:8}
Cobalah: gambarlah \omterm{def:g5:food-webs:web}{jaring makanan} sebuah \omterm{def:g2:where-animals-live:habitat}{habitat} yang kamu kenal ---
kebun, taman, kolam --- dengan sedikitnya enam \omterm{def:g5:biodiversity:species}{spesies} dan setiap anak
panah yang dapat kamu pertahankan alasannya. Tandai \omterm{def:g5:food-webs:roles}{produsennya} dengan warna
hijau.
\end{exercise}

\begin{exercise}[$\star\star$]\label{exo:g5:food-webs:9}
Jalankan \cref{met:g5:food-webs:predict} atas penyingkiran kadal dari
gambarnya: akibat langkah pertama, lalu satu langkah lagi.
\end{exercise}

\begin{exercise}[$\star\star$]\label{exo:g5:food-webs:10}
Ceritakan kembali \cref{ex:g5:food-webs:mistake} sebagai sebuah rantai
ramalan: elangnya disingkirkan, lalu apa, lalu apa? Langkah mana yang
gagal diramalkan para pengampanyenya?
\end{exercise}

\begin{exercise}[$\star\star\star$]\label{exo:g5:food-webs:11}
Pada jaring di gambarnya, bandingkan penyingkiran pohon ek (sebuah
\omterm{def:g5:food-webs:roles}{produsen}) dengan penyingkiran alap-alap (seorang pemburu puncak).
Tetangga mana yang lebih dahulu merasakan tiap kehilangannya, ke arah
mana tiap gangguannya merambat, dan kehilangan mana yang lebih mudah
diserap jaringnya di sini? Pakai
\cref{ex:g3:food-chains:dry,ex:g3:food-chains:hunter}.
\end{exercise}
